\documentclass[10pt,landscape]{article}

\usepackage[margin=0.32in]{geometry}
\usepackage[T1]{fontenc}
\usepackage[utf8]{inputenc}
\usepackage{lmodern}
\usepackage{textcomp}
\usepackage{microtype}
\usepackage[table,dvipsnames]{xcolor}
\usepackage{array}
\usepackage{booktabs}
\usepackage{enumitem}
\usepackage{hyperref}
\usepackage{listings}
\usepackage{multicol}
\usepackage{tabularx}
\usepackage{titlesec}

\pagestyle{empty}
\setlength{\parindent}{0pt}
\setlength{\parskip}{1.5pt}
\setlength{\columnsep}{13pt}
\setlength{\columnseprule}{0.2pt}
\setcounter{secnumdepth}{0}

\definecolor{ToshBlue}{HTML}{0F4C81}
\definecolor{ToshInk}{HTML}{1F2937}
\definecolor{ToshMuted}{HTML}{5B6470}
\definecolor{ToshPanel}{HTML}{F4F7FB}
\definecolor{ToshRule}{HTML}{CBD5E1}
\definecolor{ToshCodeBg}{HTML}{F8FAFC}

\color{ToshInk}

\titleformat{\section}
  {\large\bfseries\color{ToshBlue}}
  {}
  {0pt}
  {}
\titlespacing*{\section}{0pt}{0.4em}{0.12em}

\setlist[itemize]{leftmargin=1.05em,itemsep=0.5pt,topsep=1.5pt,parsep=0pt,partopsep=0pt}

\lstdefinestyle{tosh}{
  basicstyle=\ttfamily\scriptsize,
  backgroundcolor=\color{ToshCodeBg},
  frame=single,
  rulecolor=\color{ToshRule},
  framerule=0.4pt,
  xleftmargin=0pt,
  xrightmargin=0pt,
  aboveskip=0.18em,
  belowskip=0.24em,
  columns=fullflexible,
  keepspaces=true,
  showstringspaces=false,
  breaklines=true,
  literate={°}{{\textdegree}}1
           {μ}{{\ensuremath{\mu}}}1
           {Ω}{{\ensuremath{\Omega}}}1
}
\lstset{style=tosh}

\newcommand{\sheettitle}[3]{
  {\LARGE\bfseries\color{ToshBlue} ToSh Cheat Sheet: #1\par}
  \vspace{0.08em}
  {\small\color{ToshMuted} #2 \hfill #3\par}
  \vspace{0.18em}
  {\color{ToshRule}\rule{\linewidth}{0.8pt}}
  \vspace{0.12em}
}

\newcommand{\note}[1]{
  \vspace{0.12em}
  \fcolorbox{ToshRule}{ToshPanel}{%
    \parbox{\dimexpr\linewidth-2\fboxsep-2\fboxrule\relax}{\footnotesize #1}%
  }
  \vspace{0.12em}
}

\newcommand{\tighttable}[2]{
  {\footnotesize
  \renewcommand{\arraystretch}{1.08}
  \begin{tabularx}{\linewidth}{@{}#1@{}}
  #2
  \end{tabularx}}
}


\begin{document}
\footnotesize
\sheettitle{Core Syntax}{Expressions, variables, operators, and functions}{parser/tests first}

\begin{multicols*}{3}

\section{Literals And Strings}
\begin{lstlisting}
# numbers
42  3.14  1_000_000  0xFF  0b1010  0o777

# strings
"escapes\nok"
'raw text'
"""multi-line"""
$"hello {$name}"
$'ANSI C escapes: \x1b[31m'

# specials
true  false  null
2026-03-25
127.0.0.1
\end{lstlisting}

\tighttable{>{\ttfamily\raggedright\arraybackslash}p{0.25\linewidth}X}{
\toprule
Form & Notes \\
\midrule
\texttt{"..."} & Escapes, no interpolation \\
\texttt{'...'} & Raw, no interpolation \\
\texttt{"""..."""} & Raw multi-line block \\
\texttt{\$"..."} & Interpolation with \{...\} \\
\texttt{\$'...'} & ANSI-C escape processing \\
\bottomrule
}

\section{Variables And Assignment}
\begin{lstlisting}
var name = "Alice"
var age: int = 30
$age = ($age + 1)
$name ??= "guest"
\end{lstlisting}

\begin{lstlisting}
var [first, second] = [10, 20, 30]
var { Name, Uid } = {| Name = "toast", Uid = 1000 |}

var base = {| Name = "Alice", Role = "user" |}
var admin = {| ...$base, Role = "admin" |}
\end{lstlisting}

\note{Variables must be declared with \texttt{var} before use. Assignment targets use the \texttt{\$name} form. Records use paired \texttt{\{| ... |\}} delimiters; plain braces remain blocks and destructuring patterns.}

\section{Operators}
\tighttable{>{\ttfamily\raggedright\arraybackslash}p{0.24\linewidth}X}{
\toprule
Operator & What it covers \\
\midrule
\texttt{+ - * / \%} & numeric math, date arithmetic, string/list concat \\
\texttt{== != < > <= >=} & equality and ordered comparison \\
\texttt{=~ !~} & regex match and negated regex match \\
\texttt{and or not} & boolean logic; \texttt{\&\&} and \texttt{||} also work \\
\texttt{is is-not as} & type check and cast-style conversion \\
\texttt{in not in} & membership testing \\
\texttt{?? ?. ??=} & null handling \\
\texttt{?:} & ternary conditional \\
\bottomrule
}

\begin{lstlisting}
echo ("hello" =~ "^h")
echo (42 is int)
echo ("a" in ["a", "b", "c"])
var label = $count > 0 ? "items" : "empty"
\end{lstlisting}

\section{Functions}
\begin{lstlisting}
func greet(name) {
    writeline $"Hello, {$name}!"
}

func double(x) => $x * 2
func add => $1 + $2
\end{lstlisting}

\begin{lstlisting}
func connect(host: string, port: int) { ... }
func greet(name: string, title?: string = "friend") { ... }
func print-all(items...) { for item in $items { echo $item } }
func add(a: int, b: int) -> int => ($a + $b)
\end{lstlisting}

\begin{lstlisting}
var fn = func(x) => $x * 2
invoke $fn 21
var ref = &double
[1, 2, 3] | map &double
\end{lstlisting}

\note{\textbf{See also:} the \texttt{control-flow} sheet covers \texttt{if}, loops, \texttt{try/catch/finally}, \texttt{defer}, \texttt{switch}, and \texttt{match}. A natural ToSh style is typed parameters for public helpers, block syntax for pipeline predicates, and interpolated strings for readable output.}

\end{multicols*}
\newpage

\sheettitle{Core Syntax}{Comments, ranges, nullability, spread, and callable values}{page 2 / language reference}

\begin{multicols*}{3}

\section{Comments And Identifiers}
\begin{lstlisting}
# one-line comment
echo hello  # inline comment

var my_count = 10
var total-items = 42
func process-batch() { ... }
\end{lstlisting}

\begin{itemize}
\item ToSh supports single-line comments only.
\item Hyphens are allowed in identifiers, which keeps shell-style names natural.
\item Reserved words like \texttt{func}, \texttt{class}, and \texttt{match} cannot be rebound.
\end{itemize}

\section{More String Forms}
\tighttable{>{\ttfamily\raggedright\arraybackslash}p{0.26\linewidth}X}{
\toprule
Form & Use \\
\midrule
\texttt{'''...'''} & raw multi-line block \\
\texttt{\$"""..."""} & interpolated multi-line block \\
\texttt{\$'''...'''} & ANSI-C escaped block text \\
\bottomrule
}

\begin{lstlisting}
var raw = '''line 1
line 2'''

var msg = $"""User: {$name}
Count: {$count}"""
\end{lstlisting}

\section{Ranges}
\begin{lstlisting}
1..10
1..2..10
10..1
\end{lstlisting}

\begin{lstlisting}
for i in 1..5 {
    echo $i
}
\end{lstlisting}

\begin{itemize}
\item Ranges are lazy shell values, not eagerly materialized arrays.
\item They iterate naturally in \texttt{for} loops and can be splatted into commands.
\end{itemize}

\columnbreak

\section{Null And Assignment}
\begin{lstlisting}
var name = $user ?? "Anonymous"
var len = $user?.Name?.Length ?? 0

var x = 10
$x += 5
$x *= 2
$x ??= 100
\end{lstlisting}

\tighttable{>{\ttfamily\raggedright\arraybackslash}p{0.24\linewidth}X}{
\toprule
Operator & Meaning \\
\midrule
\texttt{??} & fallback when left is null \\
\texttt{?.} & null-conditional member access \\
\texttt{??=} & assign only when target is null \\
\texttt{+= -= *= /=} & compound assignment \\
\bottomrule
}

\note{A useful parsing rule: ternary binds more loosely than \texttt{??}, so \texttt{null ?? true ? yes : no} reads like \texttt{(null ?? true) ? yes : no}.}

\section{Destructuring And Spread}
\begin{lstlisting}
var person = {| name = "Alice", age = 30 |}
var { name, age } = $person

var items = [10, 20, 30, 40]
var [first, second] = $items
var [_, _, third] = $items
\end{lstlisting}

\begin{lstlisting}
var a = [1, 2, 3]
var b = [0, ...$a, 4]

var base = {| role = "user" |}
var admin = {| ...$base, role = "admin" |}

dotnet ...$publish_args
\end{lstlisting}

\section{Computed Keys}
\begin{lstlisting}
var key = "name"
var record = {| ($key) = "Alice", age = 30 |}
echo $record.name
\end{lstlisting}

\columnbreak

\section{Callable Values}
\begin{lstlisting}
var fn = func(x) => $x * 2
invoke (func(x) => ($x * 2)) 21
var callback = &greet
[1, 2, 3] | map &double
\end{lstlisting}

\begin{lstlisting}
func t1 => test1 $1 "Jim" $2

func upper_input() {
    $tosh.Function.Input | each { _.ToUpper() }
}
\end{lstlisting}

\begin{itemize}
\item \texttt{\$1}, \texttt{\$2}, ... are positional parameters in wrapper-style arrow functions.
\item \texttt{\$tosh.Function.Input} exposes piped input inside a function body.
\item \texttt{\$tosh.Function.Args} exposes the call argument list.
\end{itemize}

\section{Member Access}
\begin{lstlisting}
"hello".ToUpper()
String.Join(" + ", ["a", "b"])
$person.Name
$person?.Manager?.Name
\end{lstlisting}

\begin{itemize}
\item Dot syntax is the preferred style for instance and static member access.
\item Use \texttt{call} or \texttt{call-method} only when the invocation has to stay dynamic.
\end{itemize}

\section{Style Hints}
\begin{itemize}
\item Use interpolated strings for user-facing output and plain quoted strings for data.
\item Use typed parameters on reusable public functions, but keep local helpers lightweight.
\item Reach for destructuring and spread when they make data reshaping shorter than repeated field-by-field assignments.
\end{itemize}

\end{multicols*}
\end{document}
